% =========================================================
% Enderton, A Mathematical Introduction to Logic, 2nd ed.
% Section 1.1 -- The Language of Sentential Logic
% Notes by William Sollers, typeset in the AIM node format.
%
% Slots marked \pending{...} are items from the review that
% are yours to write. Everything else is your notes as written,
% with wording-only fixes ("sentence" -> "sentence symbol").
% =========================================================
\documentclass[11pt]{article}
\usepackage{../aim-notes}

\title{Enderton \S1.1: The Language of Sentential Logic}
\author{William Sollers}
\date{Notes, draft 2}

\begin{document}
\maketitle

% ---------------------------------------------------------
\section{Object language and metalanguage}\label{sec:languages}

\begin{definition}[Object language]\label{def:object-language}
\source{Manin, Part I \S1.1; Tarski, Ch.~I.}
\pending{Your definition: what the object language is, what it consists of,
and that it is studied rather than used. Include an example and a
non-example.}
\end{definition}

\begin{definition}[Metalanguage]\label{def:metalanguage}
\uses{def:object-language}
\source{Manin, Part I \S1.1; Tarski, Ch.~I.}
\pending{Your definition: the language used to describe the object language,
and what lives only here --- metavariables such as $\alpha,\beta$, the
quantifiers of Chapter~1, claims such as ``$\alpha$ is a wff''. Include an
example and a non-example.}
\end{definition}

\begin{remark}[Lean]
\pending{One line: which of the two Lean plays, and which an inductive
\texttt{Formula} type plays.}
\end{remark}

% ---------------------------------------------------------
\section{Symbols}

\begin{definition}[Symbols of the sentential language]\label{def:symbols-sentential}
\source{Enderton \S1.1.}
The symbols are:
\begin{center}
\begin{tabular}{llll}
\toprule
Symbol & Name & Kind & Arity \\
\midrule
$($ & left parenthesis & punctuation (logical) & --- \\
$)$ & right parenthesis & punctuation (logical) & --- \\
$\neg$ & negation & connective (logical) & 1 \\
$\land$ & conjunction & connective (logical) & 2 \\
$\lor$ & disjunction & connective (logical) & 2 \\
$\to$ & conditional & connective (logical) & 2 \\
$\leftrightarrow$ & biconditional & connective (logical) & 2 \\
$A_1, A_2, \dots$ & sentence symbols & parameter (nonlogical) & 0 (our convention) \\
\bottomrule
\end{tabular}
\end{center}
The \emph{sentential connectives} are $\neg,\land,\lor,\to,\leftrightarrow$.
The \emph{logical symbols} are the parentheses and the connectives.
The \emph{parameters} (nonlogical symbols) are the sentence symbols. There are
infinitely many sentence symbols.
\end{definition}

\begin{quantified}
\pending{Test 7 (structure): write the signature of a sentential language,
its components and their types.}
\end{quantified}

\begin{convention}[Ours, not Enderton's]
We record sentence symbols as having arity $0$. Parentheses have no arity:
they are punctuation, not operations.
\end{convention}

\begin{assumption}[Distinctness]\label{asm:symbols-distinct}
\source{Enderton \S1.1.}
The symbols listed are pairwise distinct; for example, $A_3 \neq\ \leftrightarrow$.
\end{assumption}

\begin{assumption}[No symbol is a sequence of others]\label{asm:no-symbol-is-sequence}
\source{Enderton \S1.1.}
No symbol is a finite sequence of two or more symbols; for example,
$A_3 \neq \seq{\neg, A_4, (\,}$.
\end{assumption}

\begin{example}[Why symbols must be atomic: \texttt{A1} inside \texttt{A11}]
\label{ex:atomic-symbols}
\uses{asm:no-symbol-is-sequence}
Suppose sentence symbols were written as strings of characters, so that
$A_1$ is the string \texttt{A1} and $A_{11}$ is the string \texttt{A11}.
Then \texttt{A1} occurs inside \texttt{A11}, and replacing every occurrence of
\texttt{A1} by \texttt{B} would turn the expression \texttt{(A1 -> A11)} into
\texttt{(B -> B1)}: the replacement corrupts $A_{11}$, which never contained
$A_1$ as a symbol.

Under Assumption~\ref{asm:no-symbol-is-sequence}, $A_{11}$ is a single atomic
symbol, not a sequence with $A_1$ as a part. An occurrence of $A_1$ in an
expression is therefore an entry of the sequence equal to $A_1$, and
replacing occurrences of $A_1$ cannot touch $A_{11}$.
\end{example}

\begin{remark}
This matters for \emph{substitution} (replacing every occurrence of a
sentence symbol in a wff by a wff), which comes later. It is unrelated to
\emph{concatenation}, e.g.\ $(\alpha\to\beta)=\E{\to}(\alpha,\beta)$, which
places whole expressions into a fixed pattern and replaces nothing.
\end{remark}

\begin{proposition}[Unique decomposition of expressions]\label{prop:unique-decomposition}
\uses{asm:no-symbol-is-sequence}
\source{Enderton \S1.1.}
A finite sequence of symbols can be decomposed in only one way.
\end{proposition}

\begin{quantified}
If $\seq{a_1,\dots,a_m} = \seq{b_1,\dots,b_n}$, then $m=n$ and $a_i=b_i$ for each $i$.
\pending{Missing hypothesis from the review: add the condition on the
$a_i$ and $b_j$ without which this is false.}
\end{quantified}

\begin{question}
Why is this not immediate from equality of tuples? Look at how Enderton
defines $n$-tuples in Chapter~0.
\end{question}

% ---------------------------------------------------------
\section{Expressions}

\begin{definition}[Expression]\label{def:expression}
\uses{def:symbols-sentential}
\source{Enderton \S1.1.}
An \emph{expression} is a finite sequence of symbols. We write $\Expr$ for
the set of expressions.
\end{definition}

\begin{convention}[Notation for expressions]
We \emph{denote} an expression by writing the names of its symbols side by
side. For example, $(\neg A_1)$ denotes $\seq{(\,,\neg,A_1,\,)}$. If
$\alpha$ and $\beta$ denote expressions, $(\alpha\to\beta)$ denotes their
concatenation with the symbols shown.
\end{convention}

\begin{example}
Let $\alpha = (\neg A_1)$ and $\beta = A_2$. Then
$(\alpha\to\beta) = ((\neg A_1)\to A_2)$.
\end{example}

\begin{example}[Non-example of a formula]
The expression $((\to A_3$ is an expression but is nonsense; a logic must
exclude such expressions.
\end{example}

\begin{remark}[Translation]
An English sentence such as ``Roses are red'' is presumably true or false.
Its translation is a sentence symbol, or for a compound sentence an
expression; truth values enter only in \S1.2.
\end{remark}

% ---------------------------------------------------------
\section{Well-formed formulas, informally}

\begin{definition}[Formula-building operations]\label{def:formula-building}
\uses{def:expression}
\source{Enderton \S1.1.}
For expressions $\alpha,\beta$:
\begin{align*}
\E{\neg}(\alpha) &= (\neg\alpha), &
\E{\land}(\alpha,\beta) &= (\alpha\land\beta), &
\E{\lor}(\alpha,\beta) &= (\alpha\lor\beta),\\
\E{\to}(\alpha,\beta) &= (\alpha\to\beta), &
\E{\leftrightarrow}(\alpha,\beta) &= (\alpha\leftrightarrow\beta).
\end{align*}
\end{definition}

\begin{definition}[Well-formed formula, informal]\label{def:wff-informal}
\uses{def:formula-building}
\source{Enderton \S1.1.}
\begin{enumerate}[label=(\alph*)]
\item Every sentence symbol is a wff.
\item If $\alpha$ and $\beta$ are wffs, then so are
  \pending{the five expressions produced by the formula-building operations.}
\item No expression is a wff unless it is compelled to be one by (a) and (b).
\end{enumerate}
\end{definition}

\begin{remark}
Clause (c) is not yet rigorous. Section~\ref{sec:rigorous} makes it so.
A wff can also be pictured as a tree, e.g.\ $A_1$, $A_2$ under $\land$
giving $(A_1\land A_2)$; Enderton gives the rules for trees later (\S1.3).
\end{remark}

% ---------------------------------------------------------
\section{Making clause (c) rigorous}\label{sec:rigorous}

\begin{definition}[Construction sequence]\label{def:construction-sequence}
\uses{def:formula-building}
\source{Enderton \S1.1.}
A \emph{construction sequence} is a finite sequence
$\seq{\varepsilon_1,\dots,\varepsilon_n}$ of expressions such that for each
$i\le n$ at least one of the following holds:
\begin{enumerate}[label=(\roman*)]
\item $\varepsilon_i$ is a sentence symbol;
\item $\varepsilon_i = \E{\neg}(\varepsilon_j)$ for some $j<i$;
\item $\varepsilon_i = \E{\square}(\varepsilon_j,\varepsilon_k)$ for some
  $j<i$, $k<i$, where $\square$ is one of $\land,\lor,\to,\leftrightarrow$.
\end{enumerate}
\end{definition}

\begin{quantified}
\pending{Test 6 (predicate): ``$\seq{\varepsilon_1,\dots,\varepsilon_n}$ is a
construction sequence'' $\iff$ \dots}
\end{quantified}

\begin{remark}[Duplicates and unused entries]
A construction sequence may contain repeated entries and entries that are
not used to build the last one, so a wff has many construction sequences.
\pending{Two examples that satisfy (i)--(iii): one with a duplicate, one
with an unused entry. Check $\varepsilon_1$ first.}
\end{remark}

\begin{remark}[Our generalization]
One could allow connectives of any arity $n$, with clause (iii) reading
$\varepsilon_i=\E{\square}(\varepsilon_{j_1},\dots,\varepsilon_{j_n})$ with
every $j_r<i$. Whether an $n$-ary connective can be expressed by binary ones
is a semantic question, deferred to \S1.5.
\end{remark}

\begin{definition}[Wff, bottom-up]\label{def:wff-bottom-up}
\uses{def:construction-sequence}
\source{Enderton \S1.1.}
An expression $\alpha$ is a wff if
\pending{Test 4 (quantifier): which quantifier on construction sequences?
Not ``the''.}
\end{definition}

\begin{definition}[Closed set]\label{def:closed}
\uses{def:formula-building}
\source{Enderton \S1.1, p.~18.}
\pending{Test 6 (predicate): ``$S$ is closed under the formula-building
operations'' $\iff$ \dots\ (all five operations, negation included).}
\end{definition}

\begin{remark}[Looking ahead to \S1.4]
Enderton's \S1.1 characterizes wffs bottom-up only
(Definition~\ref{def:wff-bottom-up}). In \S1.4 he gives a second
definition ``from the top down'', as the intersection of all
\emph{inductive} sets, and proves the two agree; construction sequences are
the ``bottom up'' one. Those nodes belong to the \S1.4 notes.
\end{remark}

% ---------------------------------------------------------
\section{Induction}

\begin{theorem}[Induction principle]\label{thm:induction-principle}
\uses{def:wff-bottom-up, def:closed}
\source{Enderton \S1.1, p.~18.}
If $S$ is a set of wffs containing all the sentence symbols and closed under
the formula-building operations, then $S$ is the set of all wffs.
\end{theorem}

\begin{quantified}
\pending{Test 1: the quantifier over $S$, with its domain named.}
\end{quantified}

\begin{remark}[Proof structure]
The conclusion is a set equality, $S=\WFF$. One inclusion is free:
$S\subseteq\WFF$ is a hypothesis. So the work is $\WFF\subseteq S$: take an
arbitrary wff $\alpha$ and show $\alpha\in S$.

The only handle on an arbitrary wff is the bottom-up definition: some
construction sequence ends with $\alpha$. Walk along that sequence by strong
induction, showing every entry lies in $S$; the last entry is $\alpha$.

The induction has the shape of the definition of a construction sequence.
The first entry has nothing before it, so it can only be a sentence symbol:
sentence symbols must appear before anything is built from them. Every
later entry is either a sentence symbol, in $S$ by hypothesis, or is built
from earlier entries by a formula-building operation, in $S$ by the
inductive hypothesis and closure.
\end{remark}

\begin{proof}
By hypothesis $S\subseteq\WFF$, so it suffices to show $\WFF\subseteq S$.
Let $\alpha$ be an arbitrary wff; we show $\alpha\in S$.

Since $\alpha$ is a wff, there is a construction sequence
$\seq{\beta_1,\dots,\beta_n}$ with $\beta_n=\alpha$
(Definition~\ref{def:wff-bottom-up}). We prove by strong induction on
$k\in\{1,\dots,n\}$ that $\beta_k\in S$.

\emph{Base case, $k=1$.} There are no entries before $\beta_1$, so clauses
(ii) and (iii) of Definition~\ref{def:construction-sequence} cannot hold;
hence $\beta_1$ is a sentence symbol. Since $S$ contains every sentence
symbol, $\beta_1\in S$.

\emph{Inductive case, $k>1$.} Assume $\beta_i\in S$ for all $i<k$. By
Definition~\ref{def:construction-sequence}, at least one of the following
holds for $\beta_k$:
\begin{enumerate}[label=(\roman*)]
\item $\beta_k$ is a sentence symbol. Then $\beta_k\in S$ by hypothesis.
\item $\beta_k=\E{\neg}(\beta_j)$ for some $j<k$. By the inductive
  hypothesis $\beta_j\in S$, and $S$ is closed under $\E{\neg}$, so
  $\beta_k\in S$.
\item $\beta_k=\E{\square}(\beta_j,\beta_l)$ for some $j<k$, $l<k$ and some
  binary connective $\square\in\{\land,\lor,\to,\leftrightarrow\}$. By the
  inductive hypothesis $\beta_j,\beta_l\in S$, and $S$ is closed under
  $\E{\square}$, so $\beta_k\in S$.
\end{enumerate}
In every case $\beta_k\in S$. By strong induction, $\beta_k\in S$ for all
$k\le n$; in particular $\alpha=\beta_n\in S$.

Since $\alpha$ was arbitrary, $\WFF\subseteq S$, and with
$S\subseteq\WFF$ we conclude $S=\WFF$.
\end{proof}

\begin{remark}[Idea]
A wff is the last line of a construction sequence, and $S$ absorbs every
line in turn.
\end{remark}

\begin{definition}[Parenthesis counts]\label{def:paren-counts}
\uses{def:expression}
For an expression $\varepsilon$, let $l(\varepsilon)$ be the number of
occurrences of the symbol ``$($'' in $\varepsilon$, and $r(\varepsilon)$ the
number of occurrences of ``$)$''. In particular $l(()=1$, $r())=1$, and
$l(s)=r(s)=0$ for every other symbol $s$.
\end{definition}

\begin{remark}
$l$ and $r$ are functions $\Expr\to\mathbb{N}$, not counters: every
expression already has a value, and proofs compute those values rather than
incrementing anything.
\end{remark}

\begin{lemma}[Counts add under concatenation]\label{lem:counts-additive}
\uses{def:paren-counts}
For all expressions $\varepsilon,\delta$:
$l(\varepsilon\delta)=l(\varepsilon)+l(\delta)$ and
$r(\varepsilon\delta)=r(\varepsilon)+r(\delta)$.
\end{lemma}

\begin{proof}
\pending{Your proof (short: occurrences in $\varepsilon\delta$ are the
occurrences in $\varepsilon$ together with those in $\delta$).}
\end{proof}

\begin{remark}[Structure]
With concatenation, $\Expr$ is a monoid whose identity is the empty sequence;
it is the free monoid on the symbols. The lemma says $l$ and $r$ are monoid
homomorphisms $(\Expr,\,\cdot\,,\langle\rangle)\to(\mathbb{N},+,0)$, so each
is determined by its values on single symbols.
\end{remark}

\begin{theorem}[Unbalanced parentheses]\label{thm:unbalanced-parens}
\uses{def:wff-bottom-up, lem:counts-additive}
\source{Enderton \S1.1, p.~18 (stated there as an Example).}
Any expression with more left parentheses than right parentheses is not a wff.
\end{theorem}

\begin{quantified}
\pending{Test 5: make the hidden $\forall$ explicit.}
\end{quantified}

\begin{remark}[Proof structure]
Prove something stronger: \emph{every} wff is balanced,
$l(\alpha)=r(\alpha)$. The formation rules force this, because each one
wraps balanced inputs in exactly one ``$($'' and one ``$)$''. Make that
precise by strong induction along a construction sequence of an arbitrary
wff. The theorem is then the contrapositive.
\end{remark}

\begin{proof}
Let $\alpha$ be a wff. By Definition~\ref{def:wff-bottom-up} there is a
construction sequence $\seq{\beta_1,\dots,\beta_n}$ with $\beta_n=\alpha$.
We show by strong induction on $k\in\{1,\dots,n\}$ that
$l(\beta_k)=r(\beta_k)$.

\emph{Base case, $k=1$.} As in the proof of
Theorem~\ref{thm:induction-principle}, $\beta_1$ is a sentence symbol, so
$l(\beta_1)=0=r(\beta_1)$.

\emph{Inductive case, $k>1$.} Assume $l(\beta_i)=r(\beta_i)$ for all $i<k$.
By Definition~\ref{def:construction-sequence}, one of the following holds.
\begin{enumerate}[label=(\roman*)]
\item $\beta_k$ is a sentence symbol. Then $l(\beta_k)=0=r(\beta_k)$.
\item $\beta_k=\E{\neg}(\beta_j)=(\neg\beta_j)$ for some $j<k$. By
  Lemma~\ref{lem:counts-additive},
  \[
  l(\beta_k)=1+0+l(\beta_j)+0,\qquad r(\beta_k)=0+0+r(\beta_j)+1,
  \]
  and $l(\beta_j)=r(\beta_j)$ by the inductive hypothesis, so
  $l(\beta_k)=r(\beta_k)$.
\item $\beta_k=\E{\square}(\beta_j,\beta_m)=(\beta_j\,\square\,\beta_m)$ for
  some $j,m<k$ and binary connective $\square$. By
  Lemma~\ref{lem:counts-additive},
  \[
  l(\beta_k)=1+l(\beta_j)+0+l(\beta_m)+0,\qquad
  r(\beta_k)=0+r(\beta_j)+0+r(\beta_m)+1,
  \]
  and $l(\beta_j)=r(\beta_j)$, $l(\beta_m)=r(\beta_m)$ by the inductive
  hypothesis, so $l(\beta_k)=r(\beta_k)$.
\end{enumerate}
By strong induction $l(\beta_k)=r(\beta_k)$ for all $k\le n$; in particular
$l(\alpha)=r(\alpha)$.

So every wff has $l=r$. An expression with $l>r$ therefore is not a wff.
\end{proof}

\begin{remark}[Alternative proof via the induction principle]
\uses{thm:induction-principle}
Let $S=\{\alpha\in\WFF : l(\alpha)=r(\alpha)\}$. Then $S\subseteq\WFF$ by
definition, and every sentence symbol is in $S$. If $\alpha,\beta\in S$,
the computations in cases (ii) and (iii) above show $(\neg\alpha)\in S$ and
$(\alpha\,\square\,\beta)\in S$, so $S$ is closed under the five
operations. By Theorem~\ref{thm:induction-principle}, $S=\WFF$. The
construction sequence is not mentioned: the strong induction along it was
done once, inside the proof of the induction principle.
\end{remark}

\begin{remark}
The proof shows more than the statement: no wff has more \emph{right}
parentheses than left either. Any mismatch rules an expression out.
\end{remark}

\begin{question}
Is the converse true? If not, give a counterexample.
\end{question}

\end{document}
